\documentclass[10pt,a4paper]{amsart}
\usepackage[margin=19mm]{geometry}
\usepackage{amsmath,amssymb,mathtools}
\usepackage[scr=rsfs,cal=euler]{mathalfa}
\usepackage{xfrac}
\usepackage[unicode,hidelinks,bookmarks=false]{hyperref}
\newcommand{\ie}{i.e.\ }
\newcommand{\cf}{cf.\ }
\newcommand{\resp}{respectively\ }
\newcommand{\Subsection}{\subsection}
\providecommand{\nobreakdash}{\nobreak\hbox{-}}
\setlength{\emergencystretch}{2em}
\multlinegap0pt
\usepackage{bm}
\usepackage{bbm}
\usepackage{dsfont}
\usepackage{xspace}
\usepackage{cancel}
\usepackage{mathtools}
\usepackage{amsthm}
\makeatletter
\@ifundefined{sectioncount}{\newtheorem{sectioncount}{Sectioncount}}{}
\@ifundefined{lemma}{\newtheorem{lemma}[sectioncount]{Lemma}}{}
\makeatother
\makeatletter
\newcommand{\N}{\mathbb{N}}
\newcommand{\R}{\mathbb{R}}
\DeclareMathOperator*{\argmin}{arg\,min}
\expandafter\ifx\csname assumption\endcsname\relax
\newenvironment{assumption}[1]
{\renewcommand{\theassumptionXXX}{\textup{\textbf{(#1)}}}\begin{assumptionXXX}}
{\end{assumptionXXX}}
\fi
\newcommand{\eps}{\varepsilon}
\expandafter\ifx\csname property\endcsname\relax
\newenvironment{property}[1]
{\renewcommand{\thepropertyXXX}{\textup{\textbf{(#1)}}}\begin{propertyXXX}}
{\end{propertyXXX}}
\fi
\makeatother
\title[On the breakdown point of transport-based quantiles]{On the breakdown point of transport-based quantiles}
\author{Marco Avella-Medina, Alberto Gonz\'alez-Sanz}
\date{Source: arXiv:2410.16554v2 (CC BY 4.0)}
\begin{document}
\maketitle

\noindent\emph{Source and licence.} On the breakdown point of transport-based quantiles. Version arXiv:2410.16554v2 (CC BY 4.0), distributed under the Creative Commons Attribution 4.0 International licence, as recorded in the arXiv licence metadata for this exact version and shown on the abstract page https://arxiv.org/abs/2410.16554v2. Full rights notice: licenses/arxiv-2410.16554-CC-BY-4.0.txt.\par\smallskip

\noindent\textit{Editorial scope.} Counted unit: the Lemma whose source label is \texttt{lemma:LowerCont} in the article (source0.tex lines 831--839, source label \texttt{lemma:LowerCont}), delivered together with the article's own complete original proof of it (source0.tex lines 840--879, one \texttt{proof} environment of 40 lines). No mathematics is written, repaired or re-derived: statement and proof are the article's own text line for line. Required context retained uncounted: MonotoneQuant (lines 380--383); Lemma:TukeyPositiveCont (lines 441--452); eq:convexcone (lines 681--684); eq:cone-def-nonconvex (lines 686--690); lemma:conetocone (lines 786--794); lemma:ConvergenceToOneHalph (lines 820--826). Every definition, display and label of those lines is retained unchanged. Omitted from this package and not redistributed: 1--379, 384--440, 453--680, 685--685, 691--785, 795--819, 827--830, 880--1139. Nothing inside a retained excerpt is omitted or altered: every retained body is delivered exactly as the article's lines are written, so no omission receipt of any kind accompanies this package. Citations: the retained excerpts carry no citation command at all, so no bibliography entry of the article is transcribed and no reference is redistributed. Reference adaptation: none was needed and none was made; every reference inside the delivered unit resolves inside it, so no unresolved reference or citation is produced. Printed slips kept literal: no printed word, symbol, hypothesis or number of the counted statement or of its proof was altered. Layout and class: the article's own class is \texttt{article} with packages including tikz, inputenc, fontenc, graphicx, xcolor, babel, hyperref, fancyhdr, amsmath, amstext, amsopn, amssymb, amsthm, thmtools; the wrapper is the lane's pinned \texttt{amsart} editorial wrapper at 10pt on A4, which restates the theorem-like environments the retained text needs and adds the article's own macro definitions that the retained text needs (\texttt{\textbackslash N}, \texttt{\textbackslash R}, \texttt{\textbackslash argmin}, \texttt{\textbackslash eps}), with extra wrapper packages (bm, bbm, dsfont, xspace, cancel, mathtools). The delivered numbering is the unit's own. Assets: no retained span contains a figure environment, a graphics-inclusion command, a PDF-page inclusion, a table or a TikZ picture; no image file is copied into this package, the package declares no asset dependency and no image file is redistributed with it. Licence: version arXiv:2410.16554v2 of this article is distributed under the Creative Commons Attribution 4.0 International licence, as recorded in the arXiv licence metadata for this exact version and shown on the abstract page https://arxiv.org/abs/2410.16554v2. A full rights notice is delivered at licenses/arxiv-2410.16554-CC-BY-4.0.txt. Characters: every retained excerpt is pure ASCII, so the delivered engine, XeLaTeX with its default Latin Modern text font, reports no missing character.
\begin{equation}
    \label{MonotoneQuant}
    \langle x_1 , u_1 - u_0 \rangle +\langle x_2 , u_2- u_1 \rangle +\dots+\langle x_{m-1}, u_m - u_{m-1} \rangle+ \langle x_m, u_0 - u_m \rangle \geq 0 ,
\end{equation}

\begin{lemma}
\label{Lemma:TukeyPositiveCont} 
    Set $\mu,P\in \mathcal{P}^{a.c.}(\R^d)$. Let $\partial \varphi$ be any maximal {  cyclical} monotone extension of  the optimal transport map $\nabla \varphi$ from $\mu$ to $P$.  Then it holds that 
    \begin{equation}
        \label{eq:TukeyPositiveCont}
        \{u\in \R^d: 
    {\rm TD}(u;\mu)>0\}\subset {\rm int}\left({\rm dom}(\partial \varphi)\right).
    \end{equation}
    As a consequence, for any compact set $K\subset \{u\in \R^d: 
    {\rm TD}(u;\mu)>0\}$ it holds that 
    $\partial \varphi(K)=\bigcup_{u\in K} \partial \varphi(u)$ is  compact and nonempty.
\end{lemma}

\begin{equation}
    \label{eq:convexcone}
    \mathcal{C}_{x,e, \theta} =\{ x+v\in \R^d: \ \langle v, e\rangle \geq \cos(\theta) \|v\|\} \quad  {{\rm for}\  x\in \R^d, \ e\in \mathcal{S}^{d-1}\ {\rm and}\  \theta\in \left[-\frac{\pi}{2}, \frac{\pi}{2}\right],} 
\end{equation}

\begin{equation}
    \label{eq:cone-def-nonconvex}
    \mathcal{A}_{u,e, \theta} =\left\{y\in \R^d: \langle u-y, e\rangle\leq 
 \| u-y\|  {|\sin(\theta)|}   \right\}. 
\end{equation}

\begin{lemma}\label{lemma:conetocone}
     Let $T:{\rm dom}(T)\subset \R^d\to 2^{\R^d}$ be a monotone operator, i.e., 
     \begin{equation}
         \label{monoton}
         \langle x_1-x_2, y_1-y_2\rangle\geq 0, \quad \text{for all } x_1,x_2\in {\rm dom}(T),\ y_1\in T(x_1) \ \text{and}\  y_2\in T(x_2) .
     \end{equation}
    Set $ x\in  {\rm dom}(T) $ and  {recall the definitions \eqref{eq:convexcone} and \eqref{eq:cone-def-nonconvex}.} Then, for any  $\theta\in (0, \pi/2)$ and any $e\in \mathcal{S}^{d-1}$, it holds that $$T(\mathcal{C}_{x,e, \theta}):=\bigcup_{z\in\mathcal{C}_{x,e, \theta} } T(z)\subset \mathcal{A}_{y^x,e, \theta},$$
    for all $y^x\in T(x)$. 
 \end{lemma}

\begin{lemma}\label{lemma:ConvergenceToOneHalph}
  Set $\mu\in \mathcal{P}(\R^d)$.  Let $\{\theta_k\}_{k\in \N}\subset (0,\pi/2)$ and $\{v_k\}_{k\in \N}\subset \mathbb{R}^d$ be such that  $ \theta_k\to 0$ and $ v_k\to v$, respectively, as $k\to \infty$. Then
$$ \limsup_{k\to \infty} \mu(\mathcal{A}_{u,v_k, \theta_k}) \leq \mu(\{ y: 
      \langle  y-u, v\rangle\geq 
0  \} ), $$
 {where $\mathcal{A}_{u,v_k, \theta_k}$ is defined in \eqref{eq:cone-def-nonconvex}.}
\end{lemma}

\begin{lemma}\label{lemma:LowerCont}
Let $\mu,P\in \mathcal{P}^{a.c.}(\R^d)$ and $Q_k$ be the Lebesgue uniform distribution on the ball $k v_u+\mathbb{B}$, $k\in \N$, for some 
 $$v_u\in \argmin_{v\in \mathcal{S}^{d-1}} \mu(\{ y: 
      \langle  y-u, v\rangle\geq 
0  \} ).$$
Then, for every $\eps>{\rm TD}(u;\mu)>0$, the sequence $\{\partial \varphi_{\eps,k}(u)\}_{k\in \N}$, with  
 $\partial \varphi_{\eps,k} $  being any maximal {  cyclical} monotone extension of $\nabla \varphi_{\mu\to P+\eps (Q_k-P)}$, escapes to the horizon as $k\to +\infty$, i.e.,   
    $$\inf_{y\in \partial \varphi_{\eps,k}(u)}\|y\|\to +\infty.$$
\end{lemma}

\begin{proof}
%
Note first that as ${\rm TD}(u;\mu)>0$, \autoref{Lemma:TukeyPositiveCont} implies that $ \partial\varphi_{\eps,k}(u) \neq \emptyset$ for all $k\in \N$ and  any maximal {  cyclical} monotone extension of $\nabla \varphi_{\mu\to P+\eps (Q_k-P)}$. 
We will show that for any sequence $
\{x_k\}_{k\in \N}$, with 
$ x_k \in \partial\varphi_{\eps,k}(u)$
for all $k\in \N$, it must hold that $ \|x_k\| \to \infty.$ 
To prove the claim we argue by contradiction. To do so, we can assume, after taking subsequences, that $x_k$ belongs to some ball $R\,\mathbb{B}$ for all $k\in \N$ and some finite $R>0$. % Fix $e=\frac{x_n-K_n v}{\|x_n-K_nv\|}$ and a
Assume that $k\geq  2 (R+1) $ and let 
$ \mathcal{D}_k $ be the set of all rays arising from $ x_k $ and  {intersecting} $k v_u+\mathbb{B}$. Note that for all $k$ large enough there exists $e_k\in \mathcal{S}^{d-1}$ and $\theta_k\in (0,\pi/2)$ such that  $ {{\rm cl}}(\mathcal{D}_k) =\mathcal{C}_{x_k,e_k,\theta_k}$ as defined  in \eqref{eq:convexcone}.

On the one hand, it holds that  
\begin{equation}
    \label{eq:ToContradictlowerbound1}
    (P+\eps (Q_{k}-P))(\mathcal{C}_{x_k,e_k,\theta_k})\geq  \eps Q_{k}(\mathcal{C}_{x_k,e_k,\theta_k})= \eps  {\geq } {\rm TD}(u;\mu)+\alpha
\end{equation}
for some $\alpha>0$. On the other hand, \autoref{lemma:conetocone}  and \eqref{MonotoneQuant} imply that 
$$ (\partial \varphi_{\eps,k})^{-1}(\mathcal{C}_{x_k,e_k,\theta_k})\subset \mathcal{A}_{u,e_k,\theta_k}, $$
which yields
\begin{equation}
    \label{eq:ToContradictlowerbound2}
    \mu((\partial \varphi_{\eps,k})^{-1}(\mathcal{C}_{x_k,e_k,\theta_k})) {\leq } \mu(\mathcal{A}_{u,e_k,\theta_k}).
\end{equation}

Note that $\theta_k\to 0  $ and $e_k\to v_u$, due to $\|x_k\|\leq R$ and the fact that $k v_u+\mathbb{B}$ escapes to the horizon.   The contradiction will be found by showing that as $k\to \infty$, the optimal transport map $\nabla \varphi_{\eps,k}^*=(\nabla \varphi_{\eps,k})^{-1}$ violates the push forward condition. 
Since  $\partial  \varphi_{\eps,k}^*= (\partial  \varphi_{\eps,k})^{-1}= \{\nabla  \varphi_{\eps,k}^*(x)\}$ for $(P+\eps (Q_{k}-P))$-a.e.\ $x$ and the Borel map $\nabla  \varphi_{\eps,k}^*$ pushes $P+\eps (Q_{k}-P)$ forward to $\mu$, it must hold that 
$$ \mu(\partial  \varphi_{\eps,k}^*(\mathcal{C}_{x_k,e_k,\theta_n}))= (P+\eps (Q_{k}-P))(\mathcal{C}_{x_k,e_k,\theta_k}).$$
We use \eqref{eq:ToContradictlowerbound1} and \eqref{eq:ToContradictlowerbound2} to get 
$ \mu(\mathcal{A}_{u,e_k,\theta_k}) {\geq } {\rm TD}(u;\mu)+\alpha.$ Since $\theta_k\to 0$ and $e_k\to v_u$, \autoref{lemma:ConvergenceToOneHalph} leads to the contradiction
\begin{align*}
      {\rm TD}(u;\mu)&=\min_{v\in \mathcal{S}^{d-1}}\mu(\{ y: 
      \langle  y-u, v\rangle\geq 
0  \} )\\
&=\mu(\{ y: 
      \langle  y-u, v_u\rangle\geq 
0  \} )
\geq  {\limsup_{k\to \infty}}\mu(\mathcal{A}_{u,e_k,\theta_k}) {\geq }{\rm TD}(u;\mu)+\alpha.
\end{align*}
Therefore, the result follows. 
\end{proof}

\end{document}
